\begin{description}
\item[1行目] データファイルを読み込み，\verb|dat|オブジェクトに格納します。
\item[2-9行目] 必要なデータだけに絞り込む操作です。流れを解説しますが，他のやり方でもいいですしできあがったものが何かだけわかれば結構です。
\begin{description}
	\item[3行目] 2021年のデータだけに絞り込みます。
	\item[4行目] ネタ順に並び替えています。
	\item[5行目] 年代変数とネタ順変数はもういらないので削除してしまっています。
	\item[6行目] ロング型に変換しています。これで演者-審査員-採点のデータセットができます。
	\item[7行目] 欠損値を除外しています。実はこれがこの操作の目的で，というのも過去の審査データも入っているものですから，過去の審査員も大量に変数として含まれていて，それらが欠損値になってしまっていたのです。
	\item[8-10行目] 元のワイド型に戻しています。
	\item[11行目] 以下の行列処理のため，\verb|data.frame|型から\verb|matrix|型に変換しています。
\end{description}
